% (C) 2026 Ralf Stephan, in collaboration with Claude Code.
% Released under CC0 1.0 Universal (public-domain dedication).
% See https://creativecommons.org/publicdomain/zero/1.0/

\documentclass[11pt]{article}
\usepackage[T1]{fontenc}
\usepackage[utf8]{inputenc}
\usepackage{lmodern}
\usepackage[margin=2.5cm]{geometry}
\usepackage{amsmath,amssymb}
\usepackage{booktabs,tabularx}
\usepackage{braket}
\usepackage[hidelinks]{hyperref}

% Lean names and paths: verbatim, breakable at . _ /
\DeclareUrlCommand\lean{\urlstyle{tt}}
\def\UrlBreaks{\do\.\do\_\do\/}
\newcolumntype{L}{>{\raggedright\arraybackslash}X}
\newcolumntype{B}{>{\raggedright\arraybackslash}p{0.42\textwidth}}
\emergencystretch=6em

\title{Lean formalization of an Assumptions of Physics technical brief:\\ log}
\author{}
\date{}

\begin{document}
\maketitle

\noindent
This log records which statements of the technical brief in this directory were checked in Lean,
where the formal version differs from the brief, and any problems found in the brief.
The brief is available at \url{https://assumptionsofphysics.org/resources/briefs}.

\section*{Setup}

\begin{itemize}
\item Lean \texttt{v4.36.0-rc2}; Mathlib at commit \texttt{1b1ccf4b}; physlib (\texttt{master})
  and TauCeti (commit \texttt{e7e7f52}, on the same Mathlib) as further dependencies.
\item Build: \texttt{lake build} in the project root.
\item Every file is checked with \texttt{\#print axioms} on its main results. A result is
  recorded here only if it depends on no axioms beyond \lean{propext},
  \lean{Classical.choice} and \lean{Quot.sound}, and so has no \lean{sorry}.
\item Imports are kept to the Mathlib modules actually needed.
\item Material that belongs in Mathlib is kept in \lean{ForMathlib/}, in Mathlib's directory
  layout (library \lean{ForMathlib}), and imported by the brief files.
\end{itemize}


%======================================================================

\section{Multiple equivalent characterizations of the Schr\"odinger equation}

\paragraph{Files.}\leavevmode
\begin{itemize}
\item File: \lean{briefs/SchroedingerEvolution/SchroedingerEvolution.lean}
\item Library: \lean{SchroedingerEvolution}
\item Imports:\newline \lean{ForPhyslib.QuantumMechanics.Wigner.OneParameter} and the brief
  \lean{UnitaryEntropy} (for \lean{born}, \lean{vonNeumann} and
  \lean{pure_and_born_of_preserves})
\item Wigner library, new files in \lean{ForPhyslib/QuantumMechanics/Wigner/}:
  \lean{Kadison} (mixed-state extension, the two-state bound, Parseval averaging, Kadison's
  theorem without continuity or injectivity, pure states as extreme points, the three-way
  equivalence), \lean{OneParameter} (exclusion of anti-unitaries by composition, the phase-free
  conjugation group on $\mathbb{R}$, its generator as an inner $*$-derivation, the lift to a
  unitary one-parameter group). \lean{Density} gained finite convex combinations, affine maps
  of density operators, conjugation \lean{conjOp} by a unitary or an anti-unitary, its agreement
  \lean{mul_mul_star_eq_conjOp} with $U\rho U^\dagger$ (moved there from this brief, for the brief
  ``Schr\"odinger equation is equivalent to information conservation''), and the spectral
  decomposition. \lean{OneParameter} imports physlib's
  \lean{Physlib.Mathematics.ForMathlib.OneParameterSubgroups.Unitary} (bounded Stone theorem,
  \lean{UnitaryOneParameterGroup}).
\item New file \lean{ForPhyslib/QuantumMechanics/Wigner/EqualMixtures.lean}: a map that
  preserves equal mixtures of density operators is affine on them; \lean{IsDensity.norm_le_one}
  moved here from \lean{VonNeumannResolvent}.
\item New \lean{ForMathlib} material:\newline
  \lean{ForMathlib/Analysis/Convex/Midpoint.lean} (a bounded solution of Jensen's functional
  equation on $[0,1]$ is affine, so a bounded map preserving midpoints on a convex set preserves
  convex combinations),\newline
  \lean{ForMathlib/Analysis/InnerProductSpace/Positive.lean} (a positive operator kills every
  vector on which its quadratic form vanishes),\newline
  \lean{ForMathlib/Analysis/Calculus/Deriv/ContinuousLinearMapApply.lean} (derivative of
  $s \mapsto f(s)\,x$ for operators linear over a larger field),\newline
  and in \lean{ForMathlib/Analysis/InnerProductSpace/LinearMap.lean}: a norm-preserving ring
  endomorphism of $\mathbb{C}$ is the identity or conjugation, semilinear isometries conjugate the
  inner product, and every continuous derivation of the bounded operators on an inner product
  space is inner (\lean{InnerProductSpace.exists_eq_mul_sub_mul_of_map_mul}).
\end{itemize}

\paragraph{Setting.} $E$ is a finite-dimensional complex Hilbert space, with density operators,
pure states \lean{proj}~$\psi = \ket{\psi}\bra{\psi}$, \lean{born} and \lean{vonNeumann} as in
the brief ``Schr\"odinger equation is equivalent to information conservation''. A map $T$ on
operators is \lean{AffineOnDensities} if $T(t\rho + (1-t)\sigma) = tT\rho + (1-t)T\sigma$ for
density operators $\rho, \sigma$ and $t \in [0,1]$. \lean{IsPure}~$\rho$ means $\rho = P_\psi$ for
a unit $\psi$, \lean{BijOnPure}~$T$ that $T$ maps the pure states bijectively onto the pure
states, and \lean{ImplementedByUnitaryOrAntiunitary}~$T$ that $T\rho = W\rho W^{-1}$ on density
operators for a unitary or an anti-unitary $W$. An evolution is a family $\Phi_t$ of maps on
operators, used for $t \ge 0$ and on density operators only. \lean{IsEvolution}~$\Phi$
(QST:H + [B]EVMIX + EVTI) says that each $\Phi_t$ maps density operators to density operators and
is affine on them, $\Phi_0 = \mathrm{id}$, $\Phi_{s+t} = \Phi_s \circ \Phi_t$, and
$t \mapsto \Phi_t(\rho)$ is continuous at $0^+$. The conditions are
\begin{itemize}
\item \lean{DrSceq}~$\hbar$: $\Phi_t(\rho) = U_t\rho U_t^\dagger$ with $U_0 = 1$ and
  $U'(t) = (i\hbar)^{-1}HU(t)$ for a self-adjoint $H$;
\item \lean{DrUnit}: $\Phi_t(\rho) = U_t\rho U_t^\dagger$ for physlib's
  \lean{UnitaryOneParameterGroup} $U$ (unitary, $U_0 = 1$, $U_{s+t} = U_sU_t$ on $\mathbb{R}$,
  norm-continuous);
\item \lean{DrEv}: each $\Phi_t$ is \lean{BijOnPure};
\item \lean{DrProb}: each $\Phi_t$ maps pure states to pure states (\lean{MapsPure}) and
  preserves the Born rule between them (\lean{PreservesBorn});
\item \lean{DrInfo}: each $\Phi_t$ preserves the von Neumann entropy of density operators.
\end{itemize}

\paragraph{Results.}
\begin{center}
\begin{tabularx}{\textwidth}{@{}BL@{}}
\toprule
Brief & Lean \\
\midrule
Proposition 1, complete & \lean{affine_pure_state_correspondence} (from
  \lean{Wigner.kadison_tfae}) \\
Proposition 1, (i) $\Rightarrow$ (iii) & \lean{Wigner.kadison}: onto the pure states suffices \\
Proposition 1, mixed states from pure states & \lean{Wigner.AffineOnDensities.eq_conjOp} \\
Proposition 1, (iii) $\Rightarrow$ (ii) & \lean{Wigner.bijOn_of_eq_conjOp} \\
Proposition 1, extreme points are the pure states & \lean{Wigner.isPure_iff_isExtremeDensity} \\
Proposition 1, (ii) $\Rightarrow$ (i) & \lean{Wigner.AffineOnDensities.bijOnPure} \\
Remark after Proposition 1, injective suffices & \lean{bijOnPure_of_injOn} (from
  \lean{Wigner.AffineOnDensities.bijOnPure_of_injOn}) \\
Remark, orthogonal pure states stay orthogonal &
  \lean{Wigner.AffineOnDensities.inner_eq_zero_of_injOn} \\
Proposition 1 for maps preserving equal mixtures &
  \lean{affine_pure_state_correspondence_of_map_mix} (from
  \lean{Wigner.affineOnDensities_of_map_mix}) \\
Equal mixtures preserved $\Rightarrow$ affine & \lean{Wigner.affineOnDensities_of_map_mix}, from
  \lean{Convex.map_add_of_map_midpoint} and \lean{eq_sub_smul_add_smul_of_map_midpoint} \\
Theorem 3, DR-EV $\Rightarrow$ DR-UNIT & \lean{DrEv.drUnit} \\
Theorem 3, anti-unitary case excluded by $\Phi_t = \Phi_{t/2}^2$ &
  \lean{Wigner.exists_unitary_of_map_add} \\
Theorem 3, ``the phases can be chosen'' & \lean{Wigner.exists_addChar_adL},
  \lean{Wigner.exists_unitaryOneParameterGroup_of_addChar},
  \lean{Wigner.exists_unitaryOneParameterGroup}, \lean{drUnit_of_implemented} \\
Theorem 3, DR-PROB $\Rightarrow$ DR-UNIT & \lean{DrProb.drUnit} (via
  \lean{implemented_of_mapsPure}) \\
Theorem 3, DR-INFO $\Rightarrow$ DR-UNIT & \lean{DrInfo.drProb}, then \lean{DrProb.drUnit} \\
Theorem 3, DR-UNIT $\Rightarrow$ DR-EV, DR-PROB, DR-INFO & \lean{DrUnit.drEv},
  \lean{DrUnit.drProb}, \lean{DrUnit.drInfo} \\
Theorem 3, DR-UNIT $\Leftrightarrow$ DR-SCEQ & \lean{drUnit_iff_drSceq} (for $\hbar \ne 0$),
  \lean{eq_exp_of_schrodinger} \\
Theorem 3, complete & \lean{schrodinger_tfae} (for $\hbar \ne 0$) \\
\bottomrule
\end{tabularx}
\end{center}

\paragraph{Differences from the brief.}
\begin{itemize}
\item Finite dimension only; the brief has a separable Hilbert space. In finite dimension the
  trace class is all of $\mathrm{End}(E)$, strong and norm continuity agree, and the generator is
  bounded, so Stone's theorem is physlib's bounded version.
\item The brief's evolution is defined for $t \ge 0$ only ([B]EVMIX has $t_1 \ge t_0$), so all
  conditions quantify over $t \ge 0$. DR-UNIT still asserts a unitary group on all of
  $\mathbb{R}$, as the brief's $U_{-t} = U_t^\dagger$ does.
\item Joint continuity of $(\rho, t_0, t_1) \mapsto \Phi_{t_0 \to t_1}(\rho)$ is weakened to
  continuity of $t \mapsto \Phi_t(\rho)$ at $0^+$ for each density operator $\rho$, which joint
  continuity implies. Proposition~1 needs no continuity at all. The two-parameter
  form with joint continuity is in Section~\ref{sec:schroedinger-evolution-two-parameter}.
\item DR-SCEQ is stated with a constant $H$; the brief writes $H(t)$. The brief's form is in
  Section~\ref{sec:schroedinger-evolution-two-parameter}, where it is shown to be equivalent
  under EVTI. It is stated for the propagator $U$, which gives the brief's form
  $\psi(t) = U(t)\psi(0)$ for every initial vector.
\item The ``reconstructed past'' clauses of DR-PROB and DR-INFO are dropped: a past state $\psi_P$
  is one with $\Phi_t(\psi_P) = \psi_F$, so its equality follows from the future one. In DR-EV
  they become injectivity and surjectivity on pure states.
\item DR-PROB $\Rightarrow$ DR-UNIT goes directly through Wigner's theorem in its non-bijective
  form (Geh\'er) and the mixed-state extension, not through DR-EV. DR-INFO $\Rightarrow$ DR-UNIT
  goes through DR-PROB: zero entropy characterizes pure states, and the entropy of an equal
  mixture of two pure states determines their transition probability
  (\lean{pure_and_born_of_preserves}, with equal mixtures preserved by affinity,
  \lean{Wigner.AffineOnDensities.map_mix}). Neither needs bijectivity on pure states.
\item DR-UNIT $\Rightarrow$ DR-PROB compares $\operatorname{tr}(P_{\psi'}P_{\phi'}) =
  |\braket{\psi'|\phi'}|^2$, so the phases of the image vectors never enter.
\item The proof of DR-SCEQ $\Rightarrow$ DR-UNIT shows $U_t = \exp(t(i\hbar)^{-1}H)$ because
  $\exp(-t(i\hbar)^{-1}H)U_t$ has derivative zero; it needs no $\hbar \ne 0$. The converse,
  $H = \hbar\cdot$generator, does.
\item \lean{DrSceq} and \lean{DrUnit} use the adjoint and physlib's
  \lean{UnitaryOneParameterGroup}, which require \lean{CompleteSpace E}. The brief file therefore
  carries \lean{[CompleteSpace E]} next to \texttt{[FiniteDimensional $\mathbb{C}$ E]}, although
  the second implies the first.
\end{itemize}

\paragraph{Problems found in the brief.} No errors in the statements. Gaps and implicit points:
\begin{itemize}
\item Proposition~1, (i) $\Rightarrow$ (iii), rests on ``the standard linear rank-one preserver
  result'' without a reference; it is not in any Lean library. In finite dimension it is
  replaced by a direct argument: the two-state decomposition
  $\tfrac12(P_\psi + P_\phi) = \lambda P_{\hat r} + (1-\lambda)P_{\hat s}$ with
  $\lambda = (1 + |\braket{\psi|\phi}|)/2$ shows that affinity cannot decrease
  $|\braket{\psi|\phi}|$, Parseval averaging over an orthonormal basis turns this into equality,
  and Wigner's theorem finishes. Only surjectivity onto the pure states is used; injectivity is
  not needed, and surjectivity cannot be dropped (the constant map $\rho \mapsto P_e$ is affine
  and maps pure states to pure states).
\item The Remark claims that injectivity on the pure states suffices in finite dimension ``as the
  compactness of $P(H)$ will force $T$ to be bijective''. The claim is true (for $T$ mapping pure
  states to pure states), but compactness alone does not give it: $x \mapsto x/2$ is a continuous
  injective, not surjective self-map of $[0,1]$. Invariance of domain on the manifold
  $P(\mathbb{C}^n)$ would, but it is not in any Lean library. The formalization
  (\lean{bijOnPure_of_injOn}) uses no topology. For orthogonal unit $\psi$, $\phi$ with images
  $P_a$, $P_b$, let $L$ be the linear extension of $T$ and $N = L(\ket{\phi}\bra{\psi})$. Then
  $A(z) = L(P_{\psi + z\phi}) = P_a + |z|^2P_b + zN + \bar z N^\dagger$ satisfies
  $A(z)^2 = (1 + |z|^2)A(z)$, and comparing coefficients of $z$ gives $P_aN + NP_a = N$,
  $P_bN + NP_b = N$ and $N^2 = 0$. Such an $N$ is $\ket{Na}\bra{a}$ or $\ket{a}\bra{N^\dagger a}$,
  and likewise for $b$, which forces $N = 0$ unless $a \perp b$ or $|\braket{a|b}| = 1$; and
  $N = 0$ makes $A(1) = P_a + P_b$, which is twice a pure state only if $|\braket{a|b}| = 1$.
  Injectivity excludes $|\braket{a|b}| = 1$, so orthogonality is preserved. The images of an
  orthonormal basis are then orthonormal, Bessel's inequality replaces surjectivity in the
  Parseval averaging above, and Wigner's theorem finishes.
\item ``The phases can be chosen so that $U_0 = I$ and $U_{t+s} = U_tU_s$'' is a lifting theorem
  (Bargmann's theorem for $\mathbb{R}$) that the brief does not prove. In finite dimension it is
  avoided by working with the phase-free maps $A_t = U_t(\cdot)U_t^\dagger$ on
  $\mathrm{End}(E)$: they form a norm-continuous group on $\mathbb{R}$ (the pure states span
  $\mathrm{End}(E)$), so $A_t = \exp(tL)$ in $\mathrm{End}(\mathrm{End}(E))$. $L$ is a
  $*$-derivation, every derivation of $\mathrm{End}(E)$ is inner, $L = [K, \cdot\,]$, and the
  skew part $K_0 = \tfrac12(K - K^\dagger)$ has the same commutator. Then $V_t = \exp(tK_0)$ is
  the required unitary group, and $A_t = V_t(\cdot)V_t^\dagger$.
\item The exclusion of anti-unitaries by $\Phi_t = \Phi_{t/2} \circ \Phi_{t/2}$ needs EVTI and
  the group law. In the brief ``Schr\"odinger equation is equivalent to information
  conservation'' it was not available, and continuity was used instead.
\item DR-PROB $\Rightarrow$ DR-EV: ``$p = 1$ iff same state'' gives injectivity on pure states;
  surjectivity has to come from the reconstructed past states. The direct route through Wigner's
  theorem needs neither.
\item DR-INFO $\Rightarrow$ DR-UNIT: zero entropy shows that pure states, and only pure states,
  map to pure states. Bijectivity on the pure states, which Proposition~1 needs, does not follow
  from this alone: injectivity needs the entropy of mixtures, surjectivity the reconstructed past.
  The route through DR-PROB needs neither, and no continuity is used in this step.
\item The overall phase of $U_t$ and the additive constant in $H$ (the brief's last remark) appear
  in the lift as the freedom in $K$: $K + c$ has the same commutator for every scalar $c$.
  Taking the skew part removes the real part of $c$; the imaginary part is that constant.
\end{itemize}

%======================================================================
\subsection{The conditions as the brief writes them}\label{sec:schroedinger-evolution-two-parameter}

\paragraph{Files.}\leavevmode
\begin{itemize}
\item File: \lean{briefs/SchroedingerEvolution/SchroedingerEvolution/TwoParameter.lean}
\item Library: \lean{SchroedingerEvolution}
\item Imports:\newline the finite-dimensional part \lean{SchroedingerEvolution}, which gained
  \lean{hasDerivAt_unitaryOneParameterGroup} (a unitary group solves the Schr\"odinger equation
  with $H = \hbar\cdot$generator), split off from \lean{drUnit_iff_drSceq}
\end{itemize}

\paragraph{Setting.} As in finite dimension above, with the brief's two-parameter evolution in
place of the one-parameter $\Phi_t$. A family $\Psi_{t_0 t_1} = \Phi_{t_0 \to t_1}$ of maps on
operators is used for $t_0 \le t_1$ and on density operators only.
\lean{TwoParameter.IsEvolution}~$\Psi$ (QST:H + [B]EVMIX) says that each $\Psi_{t_0 t_1}$ maps
density operators to density operators and is affine on them, $\Psi_{tt} = \mathrm{id}$,
$\Psi_{t_1 t_2} \circ \Psi_{t_0 t_1} = \Psi_{t_0 t_2}$ for $t_0 \le t_1 \le t_2$, and
$(\rho, t_0, t_1) \mapsto \Psi_{t_0 t_1}(\rho)$ is jointly continuous on the density operators and
$t_0 \le t_1$. \lean{TwoParameter.Evti}~$\Psi$~$\Phi$ (EVTI) says $\Psi_{t_0 t_1} =
\Phi_{t_1 - t_0}$ on density operators. \lean{TwoParameter.DrSceq}~$\hbar$~$\Psi$ is DR-SCEQ with a
time-dependent Hamiltonian: there are self-adjoint $H(t)$ and $U$ with $U(0) = 1$ and $U'(t) =
(i\hbar)^{-1}H(t)U(t)$ for every $t \in \mathbb{R}$, and $\Psi_{t_0 t_1}(\rho) = V\rho V^\dagger$
with the propagator $V = U(t_1)U(t_0)^\dagger$. DR-UNIT, DR-EV, DR-PROB and DR-INFO are the
conditions of the one-parameter section on $\Phi$, since the brief states them for $\Phi_t$.

\paragraph{Results.}
\begin{center}
\begin{tabularx}{\textwidth}{@{}BL@{}}
\toprule
Brief & Lean \\
\midrule
{[B]EVMIX} + EVTI give the one-parameter evolution & \lean{TwoParameter.IsEvolution.isEvolution} \\
DR-SCEQ, the solution is unitary for any self-adjoint $H(t)$ &
  \lean{TwoParameter.star_mul_self_of_hasDerivAt} \\
Theorem 3, DR-SCEQ with $H(t)$ $\Leftrightarrow$ DR-SCEQ with constant $H$ &
  \lean{TwoParameter.drSceq_iff} (for $\hbar \ne 0$) \\
Theorem 3, as stated & \lean{TwoParameter.schrodinger_tfae} (for $\hbar \ne 0$) \\
\bottomrule
\end{tabularx}
\end{center}

\paragraph{Remarks.}
\begin{itemize}
\item Joint continuity is used only at $t_0 = t_1 = 0$, along $t \mapsto (\rho, 0, t)$, to get
  the continuity of $t \mapsto \Phi_t(\rho)$ at $0^+$ that \lean{IsEvolution} asks for.
\item DR-SCEQ with $H(t)$ $\Rightarrow$ DR-UNIT: $\frac{d}{dt}(U^\dagger U) = U^\dagger
  \bigl(\overline{(i\hbar)^{-1}} + (i\hbar)^{-1}\bigr) H(t) U = 0$, so $U(t)^\dagger U(t) =
  U(0)^\dagger U(0) = 1$. This needs no continuity of $H$, and in finite dimension a left inverse
  is a right inverse, so every $U(t)$ is unitary. Then $\Phi_t = \Psi_{0t}$ is conjugation by
  $U(t)$, and \lean{drUnit_of_implemented} gives DR-UNIT.
\item The converse takes the constant $H(t) = \hbar\cdot$generator of the DR-UNIT group $V$; the
  propagator is $V_{t_1}V_{t_0}^\dagger = V_{t_1 - t_0}$ by the group law, and EVTI gives
  $\Psi_{t_0 t_1} = \Phi_{t_1 - t_0}$.
\item That $H(t)$ is then constant up to a time-dependent real multiple of the identity,
  $H(t) = H + c(t)$, is not formalized; Theorem~3 does not need it. Under EVTI the propagators
  equal $e^{(t_1 - t_0)(i\hbar)^{-1}H}$ up to a phase, which is where $c(t)$ comes from (the
  brief's remark on the phase of the unitary).
\end{itemize}

%======================================================================
\subsection{Hilbert spaces of any dimension}\label{sec:schroedinger-evolution-any-dim}

\paragraph{Files.}\leavevmode
\begin{itemize}
\item File: \lean{briefs/SchroedingerEvolution/SchroedingerEvolution/InfiniteDimensional.lean}
\item File: \lean{briefs/SchroedingerEvolution/SchroedingerEvolution/ShiftSemigroup.lean} (the
  shift semigroup, see ``Problems found in the brief''), importing \lean{InfiniteDimensional}
  and \lean{Mathlib.MeasureTheory.Function.LpSpace.DomAct.Continuous} (translation is continuous
  on $L^p$)
\item Library: \lean{SchroedingerEvolution}, which now builds every module
  \lean{SchroedingerEvolution.*}
\item Imports:\newline \lean{ForPhyslib.QuantumMechanics.Stone},\newline
  \lean{UnitaryEntropy.InfiniteDimensional} (for \lean{pure_and_born_of_preserves_vonNeumannEntropy}
  and \lean{exists_proj_of_past}),\newline
  and the finite-dimensional part \lean{SchroedingerEvolution} (for \lean{MapsPure} and
  \lean{PreservesBorn}, neither of which assumes finite dimension)
\item New files in \lean{ForPhyslib/ProbabilisticTheory/HilbertSpace/TraceClass/}:
  \lean{Compact} (Hilbert--Schmidt and trace-class operators are compact) and
  \lean{CanonicalExpansion} (a positive trace-class operator is $\sum_i \mu_i
  \ket{b_i}\bra{b_i}$ in the trace norm, with the tail estimate).
\item New file \lean{ForPhyslib/QuantumMechanics/TraceNorm.lean}:\newline the map
  \lean{AffineOnDensityOperators} and, for later use, the facts that such maps contract the trace
  norm, that conjugation by a strongly continuous family of contractions is trace-norm
  continuous, and that the pure states are closed. Only the definition is used here.
\item New file \lean{ForPhyslib/QuantumMechanics/Kadison.lean}: Kadison's theorem in any
  dimension, without continuity.
\item New file \lean{ForPhyslib/QuantumMechanics/EqualMixtures.lean}: a map that preserves equal
  mixtures of density operators is \lean{AffineOnDensityOperators}.
\item New file \lean{ForPhyslib/QuantumMechanics/Bargmann.lean}: strongly continuous unitary
  groups, and Bargmann's theorem for $\mathbb{R}$, the lift of a continuous evolution by unitary or
  anti-unitary conjugations to a strongly continuous unitary group.
\item New file \lean{ForPhyslib/QuantumMechanics/Stone.lean}: Stone's theorem in Schr\"odinger form,
  on top of PhyslibAlpha's Stone theorem
  (\lean{PhyslibAlpha/ProbabilisticTheory/HilbertSpace/Unbounded/Existence/StoneReconstruction})
  and unbounded spectral theorem.
\item New \lean{ForMathlib} material:
  \lean{ForMathlib/MeasureTheory/Integral/IntervalIntegral/LocalCocycle.lean} (a continuous local
  2-cocycle on an interval of $\mathbb{R}$ is a coboundary; the circle-valued form).
\end{itemize}

\paragraph{Setting.} $E$ is a complex Hilbert space of any dimension; separability is not used.
Density operators (\lean{IsDensityOperator}), pure states and the entropy
(\lean{vonNeumannEntropy}) are as in the section ``Hilbert spaces of any dimension'' of
\lean{briefs/UnitaryEntropy/formalization-log.tex}. A map $T$ on
operators is \lean{AffineOnDensityOperators} if $T(t\rho + (1-t)\sigma) = tT\rho + (1-t)T\sigma$
for density operators $\rho, \sigma$ and $t \in [0,1]$, and
\lean{ImplementedOnDensityOperators}~$T$ means $T\rho = W\rho W^{-1}$ on density operators for a
unitary or an anti-unitary $W$. A \lean{StronglyContinuousUnitaryGroup} is a family of unitaries
$U_t$, $t \in \mathbb{R}$, with $U_0 = 1$, $U_{s+t} = U_sU_t$ and $t \mapsto U_t x$ continuous for
every $x$. \lean{SolvesSchrodinger}~$\hbar\,H\,U$, for a partially defined operator $H$ and bounded
operators $U_t$, says that for $\psi$ in the domain of $H$ every $U_t\psi$ is in the domain and
$\frac{d}{dt}U_t\psi = (i\hbar)^{-1}HU_t\psi$. \lean{IsEvolution} is as in finite dimension, with
\lean{IsDensityOperator} and \lean{AffineOnDensityOperators}; continuity of
$t \mapsto \Phi_t(\rho)$ at $0^+$ is in operator norm. The conditions are
\begin{itemize}
\item \lean{DrSceq}~$\hbar$: $\Phi_t(\rho) = U_t\rho U_t^\dagger$ with $U_0 = 1$ and
  \lean{SolvesSchrodinger}~$\hbar\,H\,U$ for a self-adjoint $H$;
\item \lean{DrUnit}: $\Phi_t(\rho) = U_t\rho U_t^\dagger$ for a
  \lean{StronglyContinuousUnitaryGroup} $U$;
\item \lean{DrEv}: each $\Phi_t$ is \lean{BijOnPure};
\item \lean{DrProb}: each $\Phi_t$ satisfies \lean{MapsPure} and \lean{PreservesBorn}, and every
  pure state is $\Phi_t$ of a pure state (\lean{HasPurePast});
\item \lean{DrInfo}: each $\Phi_t$ preserves the von Neumann entropy, and every density operator
  is $\Phi_t$ of a density operator.
\end{itemize}

\paragraph{Results.}
\begin{center}
\begin{tabularx}{\textwidth}{@{}BL@{}}
\toprule
Brief & Lean \\
\midrule
Proposition 1, complete & \lean{affine_pure_state_correspondence} (from
  \lean{AffineOnDensityOperators.tfae}) \\
Proposition 1 for maps preserving equal mixtures &
  \lean{affine_pure_state_correspondence_of_map_mix} (from
  \lean{affineOnDensityOperators_of_map_mix}) \\
Proposition 1, orthogonal pure states stay orthogonal & \lean{AffineOnDensityOperators.inner_eq_zero} \\
Proposition 1, overlaps preserved & \lean{AffineOnDensityOperators.norm_inner_eq} \\
Proposition 1, (i) $\Rightarrow$ (iii) on pure states &
  \lean{AffineOnDensityOperators.exists_unitary_or_antiunitary} \\
Proposition 1, mixed states from pure states & \lean{AffineOnDensityOperators.eq_conjOp} \\
Proposition 1, conjugation preserves density operators & \lean{IsDensityOperator.conjOp} \\
Proposition 1, extreme points are the pure states & \lean{isPure_iff_isExtremeDensityOperator} \\
Theorem 3, DR-UNIT $\Leftrightarrow$ DR-SCEQ & \lean{drUnit_iff_drSceq} (for $\hbar \ne 0$), from
  \lean{StronglyContinuousUnitaryGroup.exists_solvesSchrodinger} and
  \lean{exists_stronglyContinuousUnitaryGroup_of_solvesSchrodinger} \\
Theorem 3, DR-UNIT $\Rightarrow$ DR-EV, DR-PROB, DR-INFO & \lean{DrUnit.drEv},
  \lean{DrUnit.drProb}, \lean{DrUnit.drInfo} \\
Theorem 3, DR-EV $\Rightarrow$ DR-UNIT & \lean{DrEv.drUnit} \\
Theorem 3, anti-unitaries excluded and ``the phases can be chosen'' &
  \lean{exists_stronglyContinuousUnitaryGroup} (Bargmann's theorem for $\mathbb{R}$) \\
Bargmann, the cocycle lemma & \lean{Real.exists_eq_exp_of_localCocycle},
  \lean{Real.exists_eq_of_continuousOn_localCocycle} \\
Theorem 3, DR-PROB $\Rightarrow$ DR-UNIT & \lean{DrProb.drEv}, then \lean{DrEv.drUnit} \\
Theorem 3, DR-INFO $\Rightarrow$ DR-UNIT & \lean{DrInfo.drProb}, then \lean{DrProb.drUnit} \\
Theorem 3, complete & \lean{schrodinger_tfae} (for $\hbar \ne 0$) \\
The past clauses of DR-PROB and DR-INFO are needed & \lean{shift_semigroup_counterexample}
  (from \lean{isEvolution_shiftEvolution}, \lean{not_drEv_shiftEvolution},
  \lean{not_drUnit_shiftEvolution}) \\
\bottomrule
\end{tabularx}
\end{center}

\paragraph{Differences from the brief.}
\begin{itemize}
\item Proposition~1 assumes affinity under all convex combinations, as in finite dimension
  (\lean{AffineOnDensityOperators}); the proof uses weights other than $\tfrac12$, namely
  $(1 + |\braket{\psi|\phi}|)/2$ in the two-state decomposition and the eigenvalues of the
  canonical expansion. A hypothesis for equal mixtures only, like the one of the brief
  ``Schr\"odinger equation is equivalent to information conservation'', is enough
  (\lean{affine_pure_state_correspondence_of_map_mix}), and needs no continuity: density
  operators have operator norm at most one, so the map is bounded on them, and a bounded map
  that preserves midpoints preserves all convex combinations. On a segment, subtracting the
  affine interpolation leaves a bounded $g$ on $[0,1]$ with $g(0) = g(1) = 0$ and
  $g(t) = \tfrac12 g(2t)$ on $[0,\tfrac12]$, $g(t) = \tfrac12 g(2t-1)$ on $[\tfrac12,1]$, so
  $\sup\|g\|$ halves and is zero. The same holds in finite dimension. The trace-norm
  contraction of affine maps (in \lean{TraceNorm}) is not used: mixed states follow from pure
  states without continuity.
\item Proposition~1, (i) $\Rightarrow$ (iii), replaces the finite-dimensional Parseval averaging,
  which needs the trace of the identity, by a two-dimensional argument. Let $r \perp s$ with
  images $a, b$ and suppose $\braket{a|b} = \delta \in (0,1)$. For the normalized $f_\pm$ of
  $a \pm b$, the number $z(x) = \braket{f_-|T(P_x) f_+}$ lies on the circle
  $|z|^2 = (1-\delta^2)/4$, and affinity sends pairs at overlap $\tfrac12$ to perpendicular
  points. The vectors $r$, $(r+s)/\sqrt2$ and $(r+is)/\sqrt2$ are pairwise at overlap $\tfrac12$,
  and three pairwise perpendicular nonzero vectors do not fit in the plane. So orthogonal pure
  states stay orthogonal, overlaps are preserved, and Wigner's theorem in any dimension with
  surjectivity gives $W$. Mixed states: cut the canonical expansion
  $\rho = \sum_i \mu_i P_{b_i}$ after weight $m$, $\rho = m\rho_F + (1-m)\sigma_F$; $T$ and
  conjugation by $W$ agree on the finite mixture $\rho_F$, so they differ on $\rho$ by at most
  $2(1-m)$ in operator norm, and $m \to 1$.
\item Continuity at $0^+$ is in operator norm, which trace-norm continuity implies, and is used
  only on pure states.
\item DR-UNIT is a strongly continuous unitary group, and DR-SCEQ allows an unbounded $H$. For
  unbounded $H$ the Schr\"odinger equation holds only for initial vectors in the domain, and
  \lean{SolvesSchrodinger} requires that the domain is invariant. DR-SCEQ asks for bounded
  operators $U_t$ with $U_0 = 1$, not for unitaries; unitarity and the group law follow.
\item The ``reconstructed past'' clauses of DR-PROB and DR-INFO are kept: every pure state has a
  pure past in DR-PROB, every density operator has a past in DR-INFO. In finite dimension they
  were dropped, but in infinite dimension they are needed (see below).
\item DR-PROB $\Rightarrow$ DR-UNIT goes through DR-EV. Born preservation gives injectivity on
  pure states (two states with the same image have transition probability $1$), and the pure past
  gives surjectivity; then Proposition~1. In finite dimension the route was Wigner's theorem
  without surjectivity, which in infinite dimension gives only an isometry.
\item The lift (\lean{exists_stronglyContinuousUnitaryGroup}) replaces the finite-dimensional
  argument through $A_t = \exp(tL)$ on $\mathrm{End}(E)$, which needs norm continuity of the
  conjugation group. Anti-unitaries are excluded by $\Phi_t = \Phi_{t/2}^2$ as before. The
  unitaries $u_t$ are then determined up to a phase: an operator commuting with every pure state
  is a scalar (\lean{exists_eq_smul_of_proj_eq}). Continuity of $t \mapsto P_{u_t\psi}$ at $0^+$
  extends to every $t \ge 0$, since $u_{t_0-h}$ is $u_h^\dagger u_{t_0}$ up to phase, and by
  polarization to $t \mapsto \ket{u_tx}\bra{u_ty}$. For a fixed unit $e$, the phase with
  $\braket{e|U_te} > 0$ makes $U_tx = \braket{U_te|e}^{-1}\ket{u_tx}\bra{u_te}\,e$ strongly
  continuous near $0$. Its multiplier $\omega(a,b) = \braket{e|U_aU_be}/\braket{e|U_{a+b}e}$,
  with $U_aU_b = \omega(a,b)U_{a+b}$, is continuous, of modulus $1$, near $1$ and a local
  2-cocycle, so $\omega(a,b) = e^{i(f(a)+f(b)-f(a+b))}$ with continuous $f$. Then
  $V_t = e^{-if(t)}U_t$ is a local homomorphism on $[0,\varepsilon]$. It extends to $t \ge 0$ by
  $V_{t/n}^n$ and to $\mathbb{R}$ by $V_{-t} = V_t^\dagger$, and the group law carries strong
  continuity from $0$ to every $t$.
\item The cocycle lemma is proved without a mollifier and without derivatives. Integrate the
  cocycle identity over $c \in [0,\varepsilon]$ at $(a,b,c)$, and over $u$ between $0$ and $b$ at
  $(a,u,\varepsilon)$. With $\Phi(x) = \int_0^\varepsilon\theta(x,c)\,dc$ and
  $K(x) = \int_0^x\theta(u,\varepsilon)\,du$ this gives
  $\varepsilon\theta(a,b) = \Phi(a) + \Phi(b) - \Phi(a+b) + K(a+b) - K(a) - K(b)$, so
  $f = (\Phi - K)/\varepsilon$. The interval may be one-sided, which the lift needs, since $U_t$ is
  given only for $t \ge 0$.
\item DR-UNIT $\Rightarrow$ DR-SCEQ takes $H = -\hbar T$ with $T$ the self-adjoint closure of the
  candidate generator, $U_t = e^{itT}$ (PhyslibAlpha's Stone theorem). DR-SCEQ $\Rightarrow$
  DR-UNIT: with $T = -\hbar^{-1}H$ and $V_t = e^{itT}$ from the spectral theorem, $U_t\psi$ and
  $V_t\psi$ solve $w' = iTw$ with the same initial value for $\psi$ in the domain, so they agree
  ($\|w_1 - w_2\|^2$ has derivative zero because $T$ is symmetric); the domain is dense, so
  $U_t = V_t$.
\item As in finite dimension, $\hbar \ne 0$ is assumed in DR-UNIT $\Leftrightarrow$ DR-SCEQ.
\end{itemize}

\paragraph{Problems found in the brief.}
\begin{itemize}
\item In infinite dimension the reconstructed past clauses of DR-PROB and DR-INFO cannot be
  dropped. On $L^2(0,\infty)$ the right shifts $(V_tf)(x) = f(x-t)$ for $x > t$ and $0$ otherwise
  form a strongly continuous semigroup of isometries. $\Phi_t(\rho) = V_t\rho V_t^\dagger$
  satisfies \lean{IsEvolution}, maps pure states to pure states, preserves the Born rule and the
  entropy, but is not a bijection of the pure states and so satisfies neither DR-EV nor DR-UNIT.
  Formalized as \lean{shift_semigroup_counterexample}. $L^2(0,\infty)$ is the closed subspace
  of $L^2(\mathbb{R})$ orthogonal to the indicators of the sets of finite measure in
  $(-\infty,0)$ (\lean{halfLine}, as the type \lean{HalfLine}), and $V_t$ is translation by
  $t \ge 0$ (\lean{rightShift}), strongly continuous because translation is continuous on
  $L^2(\mathbb{R})$. Continuity of $t \mapsto V_t\rho V_t^\dagger$ at $0^+$ in operator norm uses
  that density operators are trace class, hence compact:
  $\|V_t\rho V_t^\dagger - \rho\| \le 2\|(V_t - 1)\rho\|$, and bounded operators that tend to $0$
  pointwise tend to $0$ uniformly on the compact closure of the image of the unit ball
  (\lean{tendsto_mul_of_isCompactOperator}). For $t > 0$ the normalized indicator $e$ of
  $(0,t)$ is orthogonal to the range of $V_t$, so $\ket{e}\bra{e}$ is $\Phi_t$ of no pure state
  (\lean{not_bijOnPure_conj}).
  The brief's abstract claims the infinite-dimensional case, and its definitions do contain past
  states, so the theorem holds as stated once the past is read as part of DR-PROB and DR-INFO.
\item ``The phases can be chosen so that $U_0 = I$ and $U_{t+s} = U_tU_s$'' is Bargmann's lifting
  theorem for $\mathbb{R}$; the brief neither proves nor cites it. In infinite dimension the
  finite-dimensional trick is not available, because the conjugation group is only strongly
  continuous. The proof above is Bargmann's 1954 idea specialized to $\mathbb{R}$, with continuous
  instead of measurable multipliers.
\item Proposition~1 needs no continuity in any dimension, so the continuity in the brief's
  hypotheses is redundant there. The Remark's weakening to injectivity does not extend to infinite
  dimension: conjugation by the shift on $\ell^2(\mathbb{N})$ is affine and injective on pure
  states but not onto them (\lean{shift_preserves_not_unitary}).
\item For unbounded $H$, DR-SCEQ's ``$\psi(t) = U(t)\psi(0)$ solves the Schr\"odinger equation''
  must be restricted to $\psi(0)$ in the domain of $H$, which the brief leaves implicit. Otherwise
  no problems beyond those of the finite-dimensional part.
\end{itemize}

\end{document}
